\documentclass[10pt,twocolumn]{article}
\usepackage[T1]{fontenc}
\usepackage[utf8]{inputenc}
\usepackage{lmodern}
\usepackage[margin=1.8cm,top=2.4cm,bottom=2cm,columnsep=0.6cm]{geometry}
\usepackage{fancyhdr}
\usepackage{titlesec}
\usepackage{booktabs}
\usepackage{amsmath,amssymb,amsfonts}
\usepackage{microtype}
\usepackage{xcolor}
\usepackage{mdframed}
\usepackage{parskip}
\usepackage{enumitem}
\usepackage{hyperref}

\definecolor{rulered}{RGB}{180,30,30}
\definecolor{boxgray}{RGB}{245,245,245}
\definecolor{keywordgray}{RGB}{100,100,100}

\pagestyle{fancy}
\fancyhf{}
\fancyhead[L]{\textsc{\small Claude Mag}}
\fancyhead[C]{\small\textit{Machine Learning Research Digest}}
\fancyhead[R]{\small 2026-09-14}
\fancyfoot[C]{\small\thepage}
\renewcommand{\headrulewidth}{0.8pt}

\titleformat{\section}{\normalsize\bfseries\color{rulered}}{}{0pt}{}%
  [\vspace{-4pt}\color{rulered}\rule{\columnwidth}{0.4pt}]
\titlespacing{\section}{0pt}{8pt}{4pt}

\hypersetup{colorlinks=true,urlcolor=rulered,linkcolor=black}

\begin{document}
\setlength{\parindent}{0pt}
\setlength{\parskip}{4pt}

{\small\color{keywordgray}\textbf{Keywords:}
  topological reasoning \textbullet{}
  spatial cognition \textbullet{}
  multimodal LLM \textbullet{}
  benchmark \textbullet{}
  knot theory}\par\vspace{4pt}

{\LARGE\bfseries\raggedright
  MindTopo: Can Foundation Models
  Reason in Topological Space?}\par\vspace{4pt}

{\large\itshape\raggedright
  Best Model Scores 61.4\% vs.\ 97.9\% Human on a Benchmark
  of 11,030 Topological Intuition Instances}\par\vspace{6pt}

{\small\textbf{By} Yunfei Ge, Anbang Liu, Qineng Wang,
  Johnalbert Garnica, Jianwen Lyu, Zihan Wang, Reuben Tan,
  Jianfeng Gao, Ruohan Zhang, Yining Hong, Jiajun Wu, Manling Li
  (Northwestern University; Microsoft Research; Stanford University;
  et al.)}\par\vspace{2pt}

{\small\color{keywordgray}arXiv:2609.11900}
\par\vspace{2pt}\noindent{\color{rulered}\rule{\columnwidth}{0.6pt}}\vspace{6pt}

Spatial reasoning benchmarks for foundation models have focused almost
entirely on metric properties: distances, angles, viewpoints, and
shapes. Cognitive science identifies a more fundamental layer of
spatial understanding --- topological relations that remain invariant
under continuous deformation --- as the earliest spatial concepts
acquired by infants. MindTopo evaluates foundation models on this
layer and finds a systematic blind spot: the best multimodal LLM
scores 61.4\% on 11,030 topological instances against 97.9\% human
performance, with near-random performance on knots --- the most globally
structured topological property.

\begin{mdframed}[backgroundcolor=boxgray,linecolor=rulered,linewidth=1pt,
    innerleftmargin=6pt,innerrightmargin=6pt,
    innertopmargin=6pt,innerbottommargin=6pt]
\textbf{\small AT A GLANCE}\par\vspace{4pt}
\begin{description}[leftmargin=0pt,labelsep=4pt,itemsep=2pt,parsep=0pt,
    font=\small\bfseries]
  \item[Problem:] Foundation models are not evaluated on topological
    spatial relations despite their cognitive primacy.
  \item[Why it matters:] Topological reasoning underlies planning,
    navigation, and physical interaction; a gap here limits physical AI.
  \item[Method:] 11,030 procedurally generated instances across 5
    topological properties (continuity, separation, order, enclosure,
    knots) at reasoning and planning cognitive levels.
  \item[Result:] Best model 61.4\%; human 97.9\%; knots near random
    for all 14 tested multimodal LLMs.
  \item[Evidence:] 14 MLLMs + 3 video generative models evaluated;
    ablations with chain-of-thought, metric cue removal, difficulty scaling.
\end{description}
\end{mdframed}

\section{Topology vs.\ Metric Geometry}

Topology studies properties invariant under \emph{homeomorphisms}
(continuous deformations with continuous inverses), as opposed to
geometry, which studies properties invariant under rigid motions.
The classic example: a coffee cup and a donut are topologically
equivalent (both have one handle/hole) but geometrically different.

Cognitive developmental science (Piaget 1956; Spelke et al.\ 1995)
identifies topological relations as the first spatial concepts acquired
by infants, preceding projective and metric relations. Despite this
cognitive primacy, virtually all spatial reasoning benchmarks for
foundation models measure metric or viewpoint-dependent properties.

MindTopo addresses this gap with five topological properties:

\textbf{Continuity:} Is a path or region connected without breaks?

\textbf{Separation:} Are two regions or paths disconnected?

\textbf{Order:} What is the cyclic ordering of points along a closed curve?

\textbf{Enclosure:} Does one region contain another (Jordan curve-type)?

\textbf{Knots:} What is the global entanglement class of a closed
curve in 3D space?

\section{The Knot Problem}

Knots are the hardest property in MindTopo and the most revealing
failure. A \emph{knot} is a closed curve $\gamma: S^1 \to \mathbb{R}^3$.
Two knots are equivalent (the same knot type) iff one can be
deformed into the other without cutting --- a relation that cannot
be decomposed into any set of local features.

By the Reidemeister theorem, two knot diagrams represent equivalent
knots iff one can be obtained from the other by three local moves
(Reidemeister moves I, II, III). Checking equivalence requires
reasoning about global path structure through an unbounded sequence
of local moves. Current vision models, which process local patches
and aggregate them, have no mechanism for this kind of global
topological bookkeeping.

\textbf{Result:} All 14 MLLMs score between 24--35\% on knot
instances (4-choice random = 25\%), confirming the problem is
currently unsolvable by the tested models regardless of scale or
architecture.

\section{Benchmark Design}

MindTopo contains \textbf{11,030 instances} across \textbf{13 task types}
at two cognitive levels:

\textbf{Reasoning:} Given visual input (image or short video),
identify a topological property or predict how it changes under a
described deformation.

\textbf{Planning:} The model acts as a closed-loop agent selecting
environment actions. Task completion within $T = 20$ steps is the
primary metric.

Instances are generated algorithmically using a custom topology
simulator with difficulty parameter $d \in [1,5]$ controlling the
number of topological operations required. All instances are
procedurally generated --- not scraped --- ensuring no overlap with
web training data.

\textbf{Human baseline:} 50 participants (college students with no
topology training) evaluated a stratified 500-instance subset; mean
accuracy 97.9\%.

\section{Technical Formulation}

Let $f: X \to Y$ be a continuous map between topological spaces.
MindTopo operationalizes five invariant types:

\textbf{Separation:} Two sets $A, B \subset X$ are separated if
$\bar{A} \cap B = A \cap \bar{B} = \emptyset$.

\textbf{Enclosure:} Region $A$ encloses $B$ if $B \subset \mathrm{Int}(A)$
in the ambient space (Jordan curve theorem generalization).

\textbf{Knot equivalence:} Closed curve $\gamma: S^1 \to \mathbb{R}^3$
belongs to knot class $[K]$ if it is ambient isotopic to $K$.
MindTopo uses three-crossing and four-crossing knots as test cases.

\textbf{Order invariant:} Given $n$ points on a closed curve, their
cyclic order $\sigma \in S_n / \mathbb{Z}_n$ is invariant under
homeomorphism.

\section{Experimental Findings}

\begin{center}
\begin{tabular}{lcc}
\toprule
Model & Reasoning & Planning \\
\midrule
Human                  & 98.2\% & 97.5\% \\
Gemini 1.5 Pro (best)  & 68.3\% & 54.5\% \\
GPT-4o                 & 65.1\% & 48.9\% \\
Claude 3.5 Sonnet      & 63.7\% & 46.2\% \\
Avg.\ 14 MLLMs         & 52.4\% & 38.7\% \\
\bottomrule
\end{tabular}
\end{center}

\vspace{4pt}

Results by property reveal a clear hierarchy:

\begin{center}
\begin{tabular}{lcc}
\toprule
Property & Best Model & Human \\
\midrule
Continuity  & 79.2\% & 98.8\% \\
Separation  & 74.1\% & 98.5\% \\
Order       & 65.8\% & 97.4\% \\
Enclosure   & 61.3\% & 97.1\% \\
Knots       & 29.4\% & 97.2\% \\
\bottomrule
\end{tabular}
\end{center}

Continuity and separation (local, visually apparent properties) are
relatively easier; knots (global, non-local) are near-random.

\section{Ablations and Interpretation}

\textbf{Metric cues removed:} When images are topologically equivalent
but metrically different (differently scaled loops), model accuracy
changes by less than 2\%, confirming models do not exploit metric
shortcuts.

\textbf{Chain-of-thought prompting:} CoT improves average accuracy by
4.1 points on reasoning and 3.8 on planning --- helpful but insufficient
to close the human gap.

\textbf{Difficulty scaling:} At difficulty $d = 5$, models average
31.2\% (near random) while humans score 95.4\%, indicating complex
topological configurations are qualitatively out of distribution for
current models.

\textbf{Planning gap:} All models perform better on reasoning than
planning (average gap: 13.7 points), suggesting models recognize
topological cues without grounding them in action-relevant planning
representations.

\section*{Reference}

Yunfei Ge et al.\ \textbf{MindTopo: Can Foundation Models Reason in
Topological Space?} arXiv:2609.11900, September 2026.
\url{https://arxiv.org/abs/2609.11900}

\end{document}
